\documentclass[11pt,oneside]{report}
\usepackage{estilo}
\begin{document}
\begin{titlepage}
\thispagestyle{empty}
{\sffamily\small\color{granate} UNIVERSIDAD NACIONAL DE ENTRE RÍOS\par}
{\sffamily\small Facultad de Ciencias de la Alimentación\par}
\vspace{14mm}
{\sffamily\bfseries\fontsize{35}{40}\selectfont\color{black}AutoFinalRobo2\par}
\vspace{5mm}
{\sffamily\fontsize{19}{25}\selectfont\color{azul}Diseño y desarrollo de un\par manipulador móvil omnidireccional\par}
\vspace{5mm}
{\large Proyecto de Robótica 2\par}
\vspace{7mm}
\begin{center}
\includegraphics[width=.79\linewidth,height=.40\textheight,keepaspectratio]{robot.png}
\end{center}
\vfill
\begin{tabularx}{\linewidth}{@{}P{27mm}Y@{}}
\textbf{Integrantes} & Koba Chajchir \textperiodcentered{} Santiago Baccon\\[5pt]
\textbf{Docentes} & Nahuel Medina \textperiodcentered{} Patricio Silva\\[5pt]
\textbf{Asignatura} & Robótica 2\\[5pt]
\textbf{Ciclo lectivo} & 2026\\[5pt]
\textbf{Fecha} & 12 de noviembre de 2026
\end{tabularx}
\vspace{7mm}
{\small\color{suave}La imagen representa el diseño CAD del conjunto.\par}
\end{titlepage}

\pagenumbering{roman}
\chapter*{Resumen}
\addcontentsline{toc}{chapter}{Resumen}
AutoFinalRobo2 es un proyecto de manipulación móvil que combina una plataforma de cuatro ruedas Mecanum, un brazo antropomorfo de seis articulaciones rotacionales y una pinza simétrica. El desarrollo integra diseño mecánico, electrónica, programación embebida, sensado, comunicación inalámbrica e interfaces de operación local y remota.

La base puede representar movimientos de avance, desplazamiento lateral y giro mediante una matriz cinemática que relaciona tres velocidades del chasis con cuatro velocidades de rueda. Los encoders Hall permiten estimar desplazamientos relativos y cerrar el control de velocidad. El brazo se describe mediante una cadena de transformaciones homogéneas, cinemática directa e inversa y un Jacobiano que relaciona velocidades y fuerzas. La unión entre ambos subsistemas permite expresar la pose mundial de la herramienta y estudiar la manipulación móvil como un único sistema.

El proyecto dispone de geometría CAD del manipulador y de mediciones de la base realizadas sin carga manipulada. Los desplazamientos lineales de referencia presentan errores de llegada cercanos a 10 mm para una meta de 1 m, bajo la configuración ensayada. Los giros muestran una diferencia entre rumbo real y odométrico que requiere calibración. El análisis de masa y apoyo del conjunto identifica una limitación frontal: un factor estático de estabilidad de 1,023, inferior al criterio de diseño de 1,2. Por ello, el alcance geométrico del brazo no se identifica con una envolvente operativa segura.

El resultado del trabajo es una arquitectura integrada y un modelo matemático que conecta consigna, movimiento, medición y restricciones mecánicas. La fabricación definitiva del brazo, su potencia y bus de comunicación, la calibración de articulaciones y la caracterización de la manipulación cargada constituyen las tareas de integración que permanecen abiertas.

\textbf{Palabras clave:} robótica móvil, Mecanum, manipulador 6R, cinemática, odometría, control, estabilidad.

\tableofcontents
\clearpage
\pagenumbering{arabic}
\chapter{Proyecto y diseño del sistema}
\label{cap:proyecto}

\section{Problema y propósito}
Un robot manipulador fijo puede orientar y posicionar una herramienta dentro de un espacio limitado por su geometría. Una base móvil amplía la región de trabajo, aunque introduce incertidumbre de localización y modifica la estabilidad del conjunto. AutoFinalRobo2 aborda esa integración mediante una plataforma Mecanum y un brazo de seis ejes.

El problema consiste en convertir una consigna de desplazamiento o una pose de herramienta en acciones coordinadas, manteniendo una correspondencia clara entre geometría, unidades, sentidos de giro, mediciones y restricciones físicas. La integración debe permitir explicar por qué gira cada rueda, cómo se obtiene una posición estimada, cómo cambian las coordenadas del efector y qué condiciones limitan su operación.

El objetivo general es desarrollar una plataforma móvil con capacidad de manipulación y documentar los principios que gobiernan su diseño y funcionamiento. Los objetivos específicos son:
\begin{itemize}
\item formular la cinemática de la base y la relación entre cuentas Hall y movimiento;
\item desarrollar las transformaciones, la cinemática y el Jacobiano del brazo;
\item describir motores, transmisiones, sensores, alimentación e interfaces;
\item relacionar estabilidad, masa y cargas con las posturas de trabajo;
\item comparar consignas y mediciones para interpretar el desempeño;
\item identificar las condiciones necesarias para completar la integración física.
\end{itemize}

\section{Alcance y metodología}
El sistema contempla movilidad planar, manipulación, operación local mediante pantalla y encoder, control desde un navegador y herramientas de mantenimiento. El diseño mecánico del brazo es la configuración v17 con eslabones de PVC de 25 y 20 mm de diámetro exterior y pinza de dos mandíbulas. La plataforma se apoya únicamente sobre sus ruedas.

El desarrollo se organiza en definición de requisitos, diseño CAD, modelado matemático, implementación y caracterización. Las dimensiones nominales representan geometría de diseño; las masas calculadas dependen de densidades y de la representación de componentes. Las mediciones describen únicamente las condiciones en que fueron obtenidas. Esta distinción permite formular conclusiones proporcionadas a la evidencia disponible.

\begin{table}[!htbp]
\centering\small
\caption{Requisitos funcionales y forma de abordarlos.}
\begin{tabularx}{\linewidth}{@{}P{33mm}YY@{}}
\toprule
\textbf{Necesidad} & \textbf{Solución del proyecto} & \textbf{Condición de aceptación}\\\midrule
Movilidad planar & Cuatro ruedas Mecanum y mezcla de velocidades. & Comparar recorrido y giro con una referencia externa.\\
Posicionar herramienta & Manipulador de seis ejes y definición de TCP. & Contrastar poses y respetar límites e interferencias.\\
Sujetar objetos & Pinza de dos cremalleras accionadas por un piñón. & Caracterizar fuerza de contacto y retención.\\
Operar localmente & OLED monocroma y encoder con pulsador. & Navegación consistente y cambios confirmados.\\
Operar a distancia & Web, Wi-Fi, ESP y enlace serie al ATmega. & Consignas compatibles y estado de ejecución observable.\\
Mantener estabilidad & Restricción de postura y análisis de apoyo. & Centro de masa y momentos con margen suficiente.\\
\bottomrule
\end{tabularx}
\end{table}

\section{Arquitectura general}
El ATmega2560 concentra el control de los motores de la plataforma, adquisición de sensores e interacción con la OLED y el encoder. El ESP-12S, basado en ESP8266, aporta conectividad, servidor HTTP, WebSocket y almacenamiento LittleFS. La aplicación React y Vite se ejecuta en el navegador y presenta estado, control, visualización y herramientas de exploración.

La arquitectura se completa con las definiciones compartidas de comunicación y unidades, los modelos CAD y un cargador de mantenimiento para Windows. El software de ATmega está escrito en C con avr-libc, sin entorno Arduino ni DMA; el ESP utiliza Arduino ESP8266. La OLED es el único dispositivo I2C y su acceso se centraliza en el I2C Manager.

\begin{figure}[!htbp]
\centering
\begin{tikzpicture}[node distance=8mm and 7mm,box/.style={draw=azul!35,fill=gris,rounded corners=2mm,text width=39mm,minimum height=17mm,align=center,font=\small},arr/.style={-{Stealth[length=2mm]},thick,granate}]
\node[box] (web) {Operador y Web\\Control y visualización};
\node[box,right=of web] (esp) {ESP8266\\Wi-Fi, HTTP y WS};
\node[box,right=of esp] (avr) {ATmega2560\\Control y adquisición};
\node[box,below=of web] (tools) {CAD y herramientas\\Diseño y mantenimiento};
\node[box,below=of esp] (local) {OLED y encoder\\Interfaz local};
\node[box,below=of avr] (io) {Motores y sensores\\Plataforma física};
\draw[arr] (web)--node[above=8mm,font=\scriptsize]{WebSocket}(esp);
\draw[arr] (esp)--node[above=8mm,font=\scriptsize]{UART}(avr);
\draw[arr] (avr)--(io);
\draw[arr] (avr.south west)--(local.north east);
\end{tikzpicture}
\caption{Responsabilidades funcionales de la plataforma. La potencia y el bus STS del brazo se integran como un subsistema específico.}
\label{fig:arquitectura}
\end{figure}
\FloatBarrier

\section{Configuración mecánica}
La plataforma emplea cuatro ruedas Mecanum accionadas individualmente. El brazo incorpora seis articulaciones rotacionales: giro de base, hombro, codo, giro de muñeca, inclinación de muñeca y giro final de herramienta. La pinza añade apertura, sin aumentar los seis componentes de la pose rígida del TCP.

\figura[.69]{brazo.png}{Representación CAD del manipulador y sus interfaces principales.}{fig:brazo}

\begin{table}[!htbp]
\centering\small
\caption{Parámetros nominales de geometría y estructura.}
\begin{tabularx}{\linewidth}{@{}P{51mm}P{27mm}Y@{}}
\toprule\textbf{Magnitud} & \textbf{Valor} & \textbf{Interpretación}\\\midrule
Diámetro de rueda & 80 mm & Geometría nominal de la plataforma.\\
Suma de semidistancias & 170 mm & Parámetro de mezcla Mecanum.\\
Hombro--codo & 70 mm & Distancia geométrica entre ejes.\\
Codo--J4 & 60 mm & Tramo distal al codo.\\
J4--J5 & 55 mm & Separación hasta el centro de muñeca.\\
J5--J6 & 65 mm & Tramo distal sobre el eje de muñeca.\\
J6--TCP & 78,5 mm & Desplazamiento hasta el punto de herramienta.\\
Tubo del brazo & $\varnothing 25\times1{,}2$ mm & Diámetro interior nominal 22,6 mm.\\
Tubo del antebrazo & $\varnothing 20\times1$ mm & Diámetro interior nominal 18 mm.\\
Longitudes de corte & 42 y 28,5 mm & Piezas de tubo; no son distancias entre ejes.\\
\bottomrule
\end{tabularx}
\end{table}

La estructura se divide en montura y control fijo, base y hombro, brazo, antebrazo y giro J4, muñeca J5--J6 y herramienta. El accionamiento del codo está próximo al hombro y transmite por correa. Al acercar una masa al eje se reduce su contribución gravitatoria, aunque el cálculo completo también debe incluir poleas, soportes, tensión e inercias.

Se seleccionan cinco servos STS3250 para J1 a J5 y dos STS3032 A0 para J6 y pinza. Las reducciones incrementan el par de salida y disminuyen la velocidad. La forma CAD define la ocupación geométrica; ajustes, roscas, tolerancias, orientación de impresión y referencias comerciales deben cerrarse antes de liberar fabricación.

\subsection{Transmisiones y selección de materiales}
\begin{table}[!htbp]
\centering\small
\caption{Transmisiones nominales del manipulador.}
\begin{tabularx}{\linewidth}{@{}P{19mm}YP{28mm}@{}}
\toprule\textbf{Eje} & \textbf{Transmisión} & \textbf{Relación motor/salida}\\\midrule
J1 & Directa. & 1\\
J2 & Piñón de 18 y rueda de 72 dientes, módulo 0,8. & 4\\
J3 & Correa HP3M $210\times6$ y poleas de 20/50 dientes. & 2,5\\
J4 & Engranajes de 24/24 dientes. & 1\\
J5 y J6 & Directas. & 1\\\bottomrule
\end{tabularx}
\end{table}
Los tubos de PVC ofrecen secciones cerradas con masa reducida. Las uniones impresas resuelven alojamientos y enlaces entre piezas metálicas, rodamientos y actuadores. La resistencia de una pieza FDM depende del material efectivo, orientación, paredes, relleno, temperatura y fijación; una densidad de CAD no sustituye esas propiedades.

\subsection{Fabricación de las piezas por impresión 3D}
Fabricamos por impresión 3D FDM las piezas propias del brazo, como carcasas, alojamientos, soportes, uniones y elementos de transmisión. Los tubos de PVC, rodamientos, tornillería y actuadores se incorporan como componentes separados. El porcentaje de relleno interior, o \emph{infill}, se elige según la función mecánica de cada pieza: se reduce en las envolventes y aumenta donde deben transmitirse fuerzas y momentos.

\begin{table}[!htbp]
\centering\small
\caption{Criterio de relleno de las piezas impresas del brazo.}
\begin{tabularx}{\linewidth}{@{}P{37mm}P{20mm}Y@{}}
\toprule\textbf{Grupo de piezas} & \textbf{Relleno} & \textbf{Aplicación y motivo}\\\midrule
Carcasas y alojamientos no estructurales & 25\,\% & Todos los \emph{housings}, tapas, cubiertas y secciones no cargadas: dan forma, protegen y alojan componentes con menor consumo de material.\\
Estructura y soportes cargados & 80\,\% & Torre entre J1 y J2, soportes estructurales y uniones que transmiten momentos: se aumenta el relleno donde pasa la carga del brazo.\\
Engranajes y transmisión impresa & 100\,\% & Engranajes, piñones, cremalleras y demás piezas impresas de transmisión que requieren construcción maciza.\\\bottomrule
\end{tabularx}
\end{table}

\needspace{5\baselineskip}
La clasificación depende del recorrido de la carga. Una carcasa se imprime al 25\,\% cuando cumple una función de envolvente; si una zona también actúa como soporte y recibe el momento de una articulación, se trata como sección estructural al 80\,\%. Los engranajes se configuran macizos, con relleno nominal del 100\,\%, para evitar un interior de baja densidad bajo los dientes y en el cuerpo que transmite el par.

El momento aumenta con la fuerza y con su distancia perpendicular al eje, $M=F\,d_{\perp}$. Por eso la torre J1--J2 y sus soportes necesitan mayor relleno: reciben el efecto del peso y de las acciones del conjunto situado por encima. El relleno se combina con paredes, orientación de capas y fijaciones; por sí solo no determina una capacidad de carga. Los porcentajes indicados son el criterio de fabricación adoptado y no un resultado de resistencia medida.

Una definición de pieza identifica geometría y una ocurrencia identifica una instancia montada. Las cantidades y masas deben sumarse por ocurrencia. Un disco de paso de engranaje representa una interfaz cinemática, pero no equivale a un perfil de diente listo para fabricar.
\FloatBarrier

\chapter{Electrónica, software e instrumentación}
\label{cap:electronica}

\section{Distribución de responsabilidades}
La electrónica diferencia el control de plataforma, la conectividad y la actuación del manipulador. El ATmega2560 atiende motores DC, encoders, línea, ultrasonido, barrido del radar, batería e interfaz local. El ESP8266 sirve la aplicación Web y traduce solicitudes entre WebSocket y UART. El control del brazo debe corresponder al bus TTL serie de los servos STS y a sus tensiones de alimentación.

Los servos de bus no se controlan como servos RC por pulsos individuales desde un PCA9685. La integración requiere definir transmisión semidúplex, referencia eléctrica, alimentación, límites y manejo de fallos. La presencia de un componente en el modelo CAD representa ocupación y montaje, sin demostrar su integración funcional.

\begin{table}[!htbp]\centering\small
\caption{Asignación principal de señales en el Mega.}
\begin{tabularx}{\linewidth}{@{}P{38mm}YP{37mm}@{}}
\toprule\textbf{Elemento} & \textbf{Señales} & \textbf{Función}\\\midrule
Motor delantero derecho & D22, D24; PWM D9 & Puente H.\\
Motor delantero izquierdo & D26, D28; PWM D10 & Puente H.\\
Motor trasero derecho & D5, D6; PWM D11 & Puente H.\\
Motor trasero izquierdo & D7, D8; PWM D12 & Puente H.\\
Hall FR / FL & A8/A12; A9/A13 & Canales A/B.\\
Hall RR / RL & A10/A14; A11/A15 & Canales A/B.\\
Línea IR1 a IR5 & A4, A3, A2, A1, A0 & Lectura digital configurada.\\
Batería & A5 & Divisor resistivo.\\
Ultrasonido & TRIG D30; ECHO D31 & Tiempo de vuelo.\\
Servo de radar & D13 & Consigna de barrido.\\
OLED & SDA D20; SCL D21 & Único dispositivo I2C.\\
Encoder de menú & CLK D2; DT D3; SW D42 & Giro y pulsación.\\
Enlace ESP & TX1 D18; RX1 D19 & UART de operación.\\
Servicio ESP & D32 KEY; D33 RESET & Accionamiento mediante NPN.\\\bottomrule
\end{tabularx}
\end{table}

Los motores reciben energía desde los puentes H. Los GPIO transportan señales de control y no sustituyen las etapas de potencia. La masa lógica es común; una salida negativa de motor del puente H no se usa como referencia de tierra. El orden físico de los conectores Hall debe convertirse al orden lógico de ruedas antes de aplicar la cinemática.

\section{Alimentación y medida de batería}
La distribución debe considerar los motores de la base, lógica, radio, radar y servos del brazo. Las referencias nominales usadas en la selección son 12 V para STS3250 y 6 V para STS3032. La potencia debe dimensionarse por simultaneidad, arranque, caída de tensión y disipación. Para ramales de tensión $V_j$, corriente $I_j$ y eficiencia de conversión $\eta_j$:
\begin{align}
P_{\mathrm{carga}}&=\sum_j V_j I_j,\\
I_{\mathrm{bat}}&\simeq\frac{1}{V_{\mathrm{bat}}}\sum_j\frac{V_jI_j}{\eta_j}.
\end{align}
La autonomía aproximada se expresa como
\begin{equation}
t_{\mathrm{aut}}\simeq\frac{E_{\mathrm{útil}}}{P_{\mathrm{media}}},
\qquad E_{\mathrm{nom}}=V_{\mathrm{nom}}C_{\mathrm{Ah}}.
\end{equation}
La energía se expresa en Wh cuando la capacidad está en Ah. El valor útil depende de corriente, tensión de corte, envejecimiento y conversión; no se asigna una autonomía sin medir esas condiciones.

La batería se adquiere por A5 mediante dos resistencias nominales de $10\unit{k\Omega}$ y un capacitor de $100\unit{nF}$. Para una lectura $a$ de un ADC de $n$ bits:
\begin{equation}
V_{\mathrm{bat}}=\frac{a}{2^n-1}V_{\mathrm{ref}}
\frac{R_{\mathrm{sup}}+R_{\mathrm{inf}}}{R_{\mathrm{inf}}}.
\label{eq:bateria}
\end{equation}
Con resistencias iguales, la escala nominal es dos. La impedancia de Thévenin y la constante de tiempo son
\begin{equation}
R_{\mathrm{th}}=\frac{R_{\mathrm{sup}}R_{\mathrm{inf}}}{R_{\mathrm{sup}}+R_{\mathrm{inf}}},
\qquad \tau_{RC}=R_{\mathrm{th}}C=0{,}5\unit{ms}.
\end{equation}
El capacitor filtra ruido, pero no protege de sobretensión. El error de referencia, la tolerancia del divisor y las caídas de masa afectan la lectura. La tensión no equivale a un porcentaje de carga universal.

\section{Modelo de los motores DC}
El modelo electromecánico equivalente relaciona tensión $v_m$, corriente $i$, velocidad $\omega_m$ y par de carga $\tau_L$. Con resistencia $R_m$, inductancia $L_m$, constante de fuerza contraelectromotriz $K_e$, constante de par $K_t$, inercia $J_m$ y rozamiento viscoso $b_m$:
\begin{align}
v_m&=R_mi+L_m\dot i+K_e\omega_m,\label{eq:motor1}\\
J_m\dot\omega_m&=K_ti-b_m\omega_m-\tau_L.\label{eq:motor2}
\end{align}
Si $\vct{x}_m=[i\;\omega_m]^{\trans}$, ambas ecuaciones forman el sistema
\begin{equation}
\dot{\vct{x}}_m=
\underbrace{\begin{bmatrix}
-R_m/L_m&-K_e/L_m\\
K_t/J_m&-b_m/J_m
\end{bmatrix}}_{\mat{A}_m}\vct{x}_m+
\begin{bmatrix}1/L_m&0\\0&-1/J_m\end{bmatrix}
\begin{bmatrix}v_m\\\tau_L\end{bmatrix}.
\end{equation}
Los parámetros del motor instalado no están identificados; el modelo explica la estructura de su respuesta. En un puente H ideal, el PWM produce una tensión media $\overline v_m\simeq sDV_{\mathrm{bus}}$, con ciclo útil $0\le D\le1$ y sentido $s=\pm1$. Las pérdidas, zona muerta, carga y batería impiden interpretar PWM como velocidad medida.

Para una reducción $i_r=\omega_m/\omega_w$, la velocidad de rueda es $\omega_w=\omega_m/i_r$ y el par de salida se aproxima por $\tau_w=\eta i_r\tau_m$. Una inercia de carga referida al motor se divide por $i_r^2$.

\section{Sensores de entorno}
\subsection{Ultrasonido y barrido}
La distancia acústica $d_s$ se calcula a partir del tiempo de ida y vuelta $\Delta t$ y velocidad del sonido $c$:
\begin{equation}
d_s=\frac{c\Delta t}{2},
\qquad c\simeq331{,}3+0{,}606\,T_C\quad[\mathrm{m/s}],
\end{equation}
donde $T_C$ está en grados Celsius. A $20\,^{\circ}\mathrm C$, $c\simeq343{,}4\unit{m/s}$. Una cuantización temporal de $25\unit{\mu s}$ produce un incremento nominal de aproximadamente $4{,}29\unit{mm}$; ese incremento no representa exactitud.

Para un ángulo de barrido $\alpha$ y origen del sensor conocido, la proyección ideal en su plano es
\begin{equation}
\vct{p}_s=\begin{bmatrix}d_s\cos\alpha\\d_s\sin\alpha\\0\end{bmatrix},
\qquad \vct{p}_B=\mat{R}_{Bs}\vct{p}_s+\vct{t}_{Bs}.
\end{equation}
La consigna de servo no es una medida de su ángulo. El eco depende de superficie, inclinación, haz y tiempo de establecimiento. La muestra debe conservar validez y antigüedad. En el funcionamiento de plataforma considerado, el HC-SR04 se usa para observación y telemetría; no se le atribuye una función física de parada que no esté implementada.

\subsection{Barra de línea}
Los cinco sensores se interpretan como señales digitales en la configuración de trabajo. Si $b_i\in\{0,1\}$ representa presencia de línea y $x_i$ la posición transversal, el centroide es
\begin{equation}
e_{\ell}=\frac{\sum_{i=1}^5 b_i x_i}{\sum_{i=1}^5b_i},
\qquad \sum_i b_i>0.
\end{equation}
La implementación usa posiciones relativas $2,1,0,-1,-2$ y corrección proporcional de rumbo. Cuando no hay sensores activos o la muestra es inválida, el seguimiento se detiene. No se asigna error cero a una línea perdida.

\section{Software, comunicaciones e interfaz}
La cadena de control remoto se organiza en solicitud Web, transporte por ESP, recepción UART y ejecución supervisada por ATmega. El protocolo de enlace es ROBO2/perfil 2 con extensión EX01. Los mensajes deben conservar versión, unidades, longitud, correlación, resultado y vigencia. Una aceptación de comando confirma su admisión, mientras el fin de una misión se observa por el estado de ejecución.

El control se atiende nominalmente cada 10 ms. Una demora superior a 100 ms durante movimiento provoca corte de salidas; los incumplimientos repetidos tienen tratamiento adicional. Render de OLED, adquisición y tráfico comparten recursos, de modo que la suficiencia de memoria y tiempos debe evaluarse para cada configuración de software y carga.

La interfaz local tiene una OLED de $128\times64$ píxeles y un encoder. El corto confirma o alterna; desde el dashboard abre el menú completo. El largo se ejecuta al soltar después de al menos 800 ms: abre controles rápidos desde el dashboard y vuelve al contexto anterior desde el resto. Si existen cambios pendientes, la elección Sí/No es obligatoria.

La Web distingue datos reales, ausencia de conexión y simulación. Las secuencias locales pueden seguir ejecutándose sin una WebSocket abierta si el modo lo permite; las consignas manuales conservan su propia vigencia. El cargador Windows separa diagnóstico, operación y mantenimiento. El enlace de operación y servicio utiliza D18/D19; KEY y RESET del ESP pueden accionarse mediante NPN en D32/D33, manteniendo RESET del Mega libre en el servicio por software. La compatibilidad eléctrica debe comprobarse en el shield real.
\FloatBarrier

\chapter{Cinemática y control de la base móvil}
\label{cap:movilidad}

\section{Referencias, unidades y geometría}
Se define un marco de cuerpo $B$ con $x$ hacia adelante, $y$ a la izquierda y $z$ hacia arriba. El giro positivo $\psi$ es antihorario visto desde arriba. Las velocidades de cuerpo son
\begin{equation}
\vct{V}_B=\begin{bmatrix}v_x&v_y&\omega\end{bmatrix}^{\trans},
\qquad v_x,v_y\ [\mathrm{m/s}],\quad\omega\ [\mathrm{rad/s}].
\end{equation}
El orden lógico de ruedas es delantera izquierda (FL), delantera derecha (FR), trasera izquierda (RL) y trasera derecha (RR). Se usa radio nominal $r=0{,}040\unit m$ y suma de semidistancias $L=l_x+l_y=0{,}170\unit m$. El contacto efectivo de rodillos puede diferir de esa geometría nominal.

Para una rueda situada en $(x_i,y_i)$, la velocidad del punto de montaje inducida por un movimiento rígido es
\begin{equation}
\vct{v}_i=\begin{bmatrix}v_x-\omega y_i\\v_y+\omega x_i\end{bmatrix}.
\end{equation}
La proyección según los rodillos a 45 grados produce las combinaciones de avance, lateral y giro. La disposición física de rodillos debe corresponder a la matriz utilizada. El modelo sigue el formalismo de robots móviles omnidireccionales~\cite{omni}.

\section{Cinemática inversa Mecanum}
Sea $\vct{u}=[u_{\mathrm{FL}},u_{\mathrm{FR}},u_{\mathrm{RL}},u_{\mathrm{RR}}]^{\trans}$ el vector de velocidades periféricas de rueda. El mezclador geométrico es
\begin{equation}
\vct{u}=\mat{A}\mat{S}\vct{V}_B,
\qquad
\mat{A}=\begin{bmatrix}
1&-1&-L\\1&1&L\\1&1&-L\\1&-1&L
\end{bmatrix}.
\label{eq:mecanum}
\end{equation}
La matriz $\mat S$ adapta los sentidos del contrato lógico a los de movimiento. Para la configuración descrita:
\begin{equation}
\mat{S}=\diag(-1,-1,-1)=-\mat{I}_3.
\end{equation}
Estos signos no sustituyen las calibraciones eléctricas individuales del motor y del encoder. Primero se normalizan las cuentas y las salidas por rueda; luego se usan las matrices de cuerpo.

La velocidad angular accionada se obtiene como
\begin{equation}
\vct{\omega}_w=\frac{1}{r}\mat A\mat S\vct V_B.
\end{equation}
La distinción es esencial: $u_i$ es velocidad periférica, $\omega_{w,i}$ es velocidad angular y PWM es una acción eléctrica. No son magnitudes intercambiables.

\subsection{Ejemplo de avance}
Para $\vct V_B=[0{,}100\;0\;0]^{\trans}$, en unidades SI, se obtiene
\begin{equation}
\vct u=\begin{bmatrix}-0{,}100\\-0{,}100\\-0{,}100\\-0{,}100\end{bmatrix}\unit{m/s},
\qquad
\vct\omega_w=\begin{bmatrix}-2{,}5\\-2{,}5\\-2{,}5\\-2{,}5\end{bmatrix}\unit{rad/s}.
\end{equation}
La magnitud equivale a aproximadamente $23{,}87\unit{rpm}$. El signo negativo resulta de la base lógica elegida; la polaridad efectiva se aplica en el controlador de cada rueda.

\subsection{Límites y saturación}
Para un límite periférico común $u_{\max}$, un escalado uniforme conserva la dirección solicitada:
\begin{equation}
\gamma=\max\left(1,\max_i\frac{|u_i|}{u_{\max}}\right),
\qquad \vct u_{\mathrm{apl}}=\frac{\vct u}{\gamma}.
\end{equation}
Recortar las ruedas independientemente puede cambiar las proporciones de avance, lateral y giro. El parámetro nominal de 150 mm/s del software es una consigna límite de configuración, sin equivaler a una velocidad máxima física caracterizada.

\section{Cinemática directa y redundancia}
La reconstrucción ideal se obtiene por pseudoinversa. Como $L\ne0$:
\begin{equation}
\mat A^{\trans}\mat A=\diag(4,4,4L^2),
\qquad \mat A^+=(\mat A^{\trans}\mat A)^{-1}\mat A^{\trans}.
\end{equation}
Por tanto,
\begin{equation}
\mat A^+=\frac14\begin{bmatrix}
1&1&1&1\\-1&1&1&-1\\-1/L&1/L&-1/L&1/L
\end{bmatrix},
\qquad \vct V_B=\mat S^{-1}\mat A^+\vct u.
\label{eq:directam}
\end{equation}
En esta configuración $\mat S^{-1}=\mat S$. Cuatro observaciones de rueda determinan tres velocidades de cuerpo bajo el modelo ideal. El residuo
\begin{equation}
\vct\varepsilon=\vct u-\mat A\mat S\widehat{\vct V}_B
\end{equation}
permite observar incompatibilidades entre medidas y modelo. Un deslizamiento coherente de varias ruedas puede no producir un residuo grande, de modo que no reemplaza una referencia externa.

\section{Encoders y conversión de cuentas}
La configuración nominal utiliza $N_i=370$ cuentas por vuelta de rueda y captura x1: flanco ascendente del canal A y lectura de B para el sentido. Ese valor no se multiplica nuevamente por cuatro.

Para un incremento normalizado $\Delta n_i$ y período $\Delta t$:
\begin{align}
\Delta\phi_i&=\frac{2\pi\Delta n_i}{N_i},\\
\Delta s_i&=r_i\Delta\phi_i=\frac{\pi D_i\Delta n_i}{N_i},\\
\widehat u_i&=\frac{\Delta s_i}{\Delta t}.
\end{align}
Con $D=80\unit{mm}$, la cuantización lineal nominal es
\begin{equation}
q_s=\frac{\pi D}{N}=\frac{80\pi}{370}\simeq0{,}679\unit{mm/cuenta}.
\end{equation}
La cuantización no representa la exactitud del desplazamiento: diámetro efectivo, deformación, carga, piso y patinamiento pueden dominar el error. Una ausencia de cuentas puede indicar reposo, falta de conexión o fallo de adquisición; se necesita estado de validez para interpretarla.

\section{Odometría y escalas direccionales}
La inversa Mecanum también se aplica a incrementos de recorrido:
\begin{equation}
\Delta\vct\xi_B=\begin{bmatrix}\Delta x_B\\\Delta y_B\\\Delta\psi\end{bmatrix}
=\mat S^{-1}\mat A^+\begin{bmatrix}\Delta s_{\mathrm{FL}}\\\Delta s_{\mathrm{FR}}\\\Delta s_{\mathrm{RL}}\\\Delta s_{\mathrm{RR}}\end{bmatrix}.
\end{equation}
Se aplican factores empíricos por componente y por sentido antes de transformar al mundo:
\begin{equation}
\Delta\vct\xi_{B,c}=
\underbrace{\begin{bmatrix}k_x&0&0\\0&k_y&0\\0&0&k_\psi\end{bmatrix}}_{\mat K_d}
\Delta\vct\xi_B.
\end{equation}
El signo de cada incremento elige avance o retroceso, izquierda o derecha, y CW o CCW. Los valores de fábrica considerados son 1,000 para avance y retroceso, 0,775 para laterales y 1,028 para giros. Una configuración guardada puede contener otros valores. Estas escalas corrigen odometría; no multiplican PWM ni la meta de misión.

\subsection{Interpretación de la calibración}
Si una misión se detiene al alcanzar una meta odométrica $d_{\mathrm{obj}}$ y se recorre una distancia real $d_{\mathrm{real}}$, el ajuste es
\begin{equation}
k_{\mathrm{nuevo}}=k_{\mathrm{actual}}
\frac{d_{\mathrm{real}}}{d_{\mathrm{obj}}}.
\label{eq:ajustek}
\end{equation}
Si al contar 1000 mm laterales se recorren 775 mm reales, el factor pasa de 1 a 0,775. Para el mismo incremento de cuentas se publica menos desplazamiento y la misión requiere más cuentas para completar la meta. Si la parada la define una distancia física externa, se compara la distancia real con la odometría obtenida sobre ese mismo tramo; no se intercambian ambos protocolos.

Los factores tienen resolución de una milésima, rango 0 a 2, y cero bloquea la componente correspondiente en control cerrado. Modificar la escala exige reiniciar el marco de pose para no mezclar una trayectoria con dos calibraciones.

\subsection{Integración en el mundo}
La relación diferencial es
\begin{equation}
\begin{bmatrix}\dot x\\\dot y\\\dot\psi\end{bmatrix}=
\begin{bmatrix}\cos\psi&-\sin\psi&0\\\sin\psi&\cos\psi&0\\0&0&1\end{bmatrix}
\begin{bmatrix}v_x\\v_y\\\omega\end{bmatrix}.
\end{equation}
La integración de punto medio usa los incrementos corregidos y el rumbo medio
\begin{align}
\psi_m&=\psi_k+\frac{\Delta\psi_c}{2},\\
x_{k+1}&=x_k+\Delta x_{B,c}\cos\psi_m-\Delta y_{B,c}\sin\psi_m,\\
y_{k+1}&=y_k+\Delta x_{B,c}\sin\psi_m+\Delta y_{B,c}\cos\psi_m,\\
\psi_{k+1}&=\wrap(\psi_k+\Delta\psi_c).
\end{align}
La pose es relativa al origen de inicialización. Sin referencia externa acumula deriva y no constituye localización absoluta.

\section{Regulación de velocidad por rueda}
El control compara consigna y velocidad Hall con acumuladores independientes. La implementación usa ganancias comunes en formato Q8 y recortes numéricos:
\begin{align}
e_i[k]&=\sat_{2000}\!\left(u_{\mathrm{ref},i}[k]-\widehat u_i[k]\right),\\
E_i[k]&=\sat_{32767}\!\left(E_i[k-1]+e_i[k]\right),\\
p_i^*[k]&=\sat_{p_{\max}}\!\left(
\frac{K_pe_i[k]+K_iE_i[k]+K_d\bigl(e_i[k]-e_i[k-1]\bigr)}{256}\right).
\end{align}
Las velocidades de esta implementación se expresan en mm/s. Los valores iniciales son $K_p=128$, $K_i=8$ y $K_d=0$. El acumulador suma por ciclo sin multiplicación explícita por tiempo; sus ganancias dependen del período de servicio. Antes de escribir al puente H se aplica una rampa a la salida.

El recorte del acumulador evita desbordamiento, pero no equivale a un antiwindup completo ligado a saturación. La compensación de zona muerta y el ajuste de ganancias requieren medidas. PWM manual y prueba de rueda evitan este lazo, por lo que no tienen las mismas condiciones que una misión cerrada.

\section{Navegación a una pose planar}
Para una meta $(x_d,y_d,\psi_d)$ se calculan errores de mundo y se rotan al cuerpo:
\begin{equation}
\begin{bmatrix}e_{x,B}\\e_{y,B}\end{bmatrix}
=\begin{bmatrix}\cos\psi&\sin\psi\\-\sin\psi&\cos\psi\end{bmatrix}
\begin{bmatrix}x_d-x\\y_d-y\end{bmatrix},
\qquad e_\psi=\wrap(\psi_d-\psi).
\end{equation}
La ley nominal previa a límites es
\begin{equation}
v_x=K_{xy}e_{x,B},\qquad v_y=K_{xy}e_{y,B},\qquad
\omega=K_\psi e_\psi.
\end{equation}
Se utiliza $K_{xy}\simeq1\unit{s^{-1}}$ y ganancia de rumbo de $3\unit{s^{-1}}$ durante traslación, frente a $1\unit{s^{-1}}$ en giro puro. La implementación agrega límites, mínimos para superar zona muerta y normalización de mezcla. El giro puro anula las consignas de traslación.

La configuración de llegada de la fuente 1.28.1 usa
\begin{equation}
|e_x|\le4\unit{mm},\qquad |e_y|\le4\unit{mm},\qquad
|e_\psi|\le10\unit{mrad}\simeq0{,}57^{\circ}.
\end{equation}
Estas bandas son decisiones de software, no una precisión física demostrada. Los resultados experimentales del capítulo~\ref{cap:resultados} corresponden a 1.28.0, con bandas de 10 mm y 25 mrad.
\FloatBarrier

\chapter{Modelo cinemático del manipulador}
\label{cap:brazo}

\section{Convención de transformaciones}
Una pose rígida se representa con una rotación $\mat R_{AB}$ y una traslación $\vct t_{AB}$. Se adopta la formulación de movimientos rígidos~\cite{lynch}, con $\mat T_{AB}$ como transformación que lleva coordenadas expresadas en $B$ al marco $A$:
\begin{equation}
\widetilde{\vct p}_A=\mat T_{AB}\widetilde{\vct p}_B,
\qquad
\mat T_{AB}=\begin{bmatrix}
\mat R_{AB}&\vct t_{AB}\\\vct0^{\trans}&1
\end{bmatrix},
\qquad \widetilde{\vct p}=\begin{bmatrix}\vct p\\1\end{bmatrix}.
\end{equation}
Las rotaciones propias cumplen $\mat R^{\trans}\mat R=\mat I$ y $\det\mat R=1$. La composición y la inversa son
\begin{equation}
\mat T_{AC}=\mat T_{AB}\mat T_{BC},\qquad
\mat T_{AB}^{-1}=\begin{bmatrix}
\mat R_{AB}^{\trans}&-\mat R_{AB}^{\trans}\vct t_{AB}\\
\vct0^{\trans}&1
\end{bmatrix}.
\end{equation}
El orden importa: las matrices actúan de derecha a izquierda sobre vectores columna. Los vectores libres tienen última coordenada cero y no reciben la traslación aplicada a los puntos.

Las rotaciones elementales son
\begin{align}
\mat R_x(\theta)&=\begin{bmatrix}1&0&0\\0&\cos\theta&-\sin\theta\\0&\sin\theta&\cos\theta\end{bmatrix},\\
\mat R_y(\theta)&=\begin{bmatrix}\cos\theta&0&\sin\theta\\0&1&0\\-\sin\theta&0&\cos\theta\end{bmatrix},\\
\mat R_z(\theta)&=\begin{bmatrix}\cos\theta&-\sin\theta&0\\\sin\theta&\cos\theta&0\\0&0&1\end{bmatrix}.
\end{align}
Sus versiones homogéneas se obtienen agregando columna de traslación nula y fila $[0\;0\;0\;1]$. Una traslación local sobre $z$ es
\begin{equation}
\mat T_z(\ell)=\begin{bmatrix}1&0&0&0\\0&1&0&0\\0&0&1&\ell\\0&0&0&1\end{bmatrix}.
\end{equation}
En las ecuaciones se usan metros y radianes. Los grados se reservan para lectura de límites y consignas de interfaz.

\section{Montura y cadena de articulaciones}
El modelo usa seis ángulos geométricos $\vct q=[q_1,\ldots,q_6]^{\trans}$ con ejes locales $Z,Y,Y,Z,Y,Z$. Los parámetros equivalentes son
\begin{align}
h&=0{,}096+0{,}113=0{,}209\unit m,\\
L_1&=0{,}070\unit m,\qquad
L_2=0{,}060+0{,}055=0{,}115\unit m,\\
d&=0{,}065+0{,}0785=0{,}1435\unit m.
\end{align}
La agrupación de $L_2$ y $d$ aprovecha tramos colineales que no cambian de dirección al rotar alrededor de su propio eje. Así, $L_2$ no es una longitud de corte del tubo y $d$ no es únicamente la longitud física de la pinza.

La transformación desde la montura $M$ hasta la herramienta $E$ es
\begin{equation}
\begin{aligned}
\mat T_{ME}(\vct q)={}&\mat T_z(h)\mat R_z(q_1)\mat R_y(q_2)\mat T_z(L_1)\\
&\cdot\mat R_y(q_3)\mat T_z(L_2)\mat R_z(q_4)\mat R_y(q_5)
\mat R_z(q_6)\mat T_z(d).
\end{aligned}
\label{eq:cadena}
\end{equation}
En esta expresión las matrices de rotación se entienden embebidas en $4\times4$. La montura del brazo respecto de la base móvil se representa por
\begin{equation}
\mat T_{BM}=\begin{bmatrix}
0&1&0&0{,}027819\\-1&0&0&0\\0&0&1&0{,}079041\\0&0&0&1
\end{bmatrix}.
\label{eq:montura}
\end{equation}
La rotación equivale a $-90^{\circ}$ sobre $z$. La conversión de coordenadas CAD a la representación del robot y la transformación de montura cumplen funciones diferentes; no se suman como dos desplazamientos físicos del mismo origen. Las transformaciones de presentación del visor tampoco se aplican nuevamente a la pose física.

\subsection{Formulación Denavit--Hartenberg equivalente}
La convención DH estándar define
\begin{equation}
\mat A_i=\mat R_z(\theta_i)\mat T_z(d_i)\mat T_x(a_i)\mat R_x(\alpha_i).
\end{equation}
Con $c_\theta=\cos\theta_i$, $s_\theta=\sin\theta_i$, $c_\alpha=\cos\alpha_i$ y $s_\alpha=\sin\alpha_i$:
\begin{equation}
\mat A_i=\begin{bmatrix}
c_\theta&-s_\theta c_\alpha&s_\theta s_\alpha&a_ic_\theta\\
s_\theta&c_\theta c_\alpha&-c_\theta s_\alpha&a_is_\theta\\
0&s_\alpha&c_\alpha&d_i\\0&0&0&1
\end{bmatrix}.
\end{equation}
Una elección de marcos equivalente a la cadena~\eqref{eq:cadena} es la siguiente. Sus offsets pertenecen al modelo matemático y no son ceros electrónicos de servos.
\begin{table}[!htbp]\centering
\caption{Parámetros DH estándar del modelo equivalente.}
\begin{tabular}{ccccc}\toprule
$i$&$\theta_i$&$d_i\,[\mathrm m]$&$a_i\,[\mathrm m]$&$\alpha_i$\\\midrule
1&$q_1$&$h$&0&$\pi/2$\\
2&$\pi/2-q_2$&0&$L_1$&0\\
3&$\pi/2-q_3$&0&0&$\pi/2$\\
4&$q_4-\pi$&$L_2$&0&$-\pi/2$\\
5&$q_5$&0&0&$\pi/2$\\
6&$q_6$&$d$&0&0\\\bottomrule
\end{tabular}
\end{table}
El producto ordenado $\mat A_1\cdots\mat A_6$ expresa la pose final con esos marcos. La cadena por movimientos elementales se mantiene como referencia de interpretación porque muestra directamente los ejes físicos y las agrupaciones de longitud.

\section{Cinemática directa}
Se define el centro de muñeca $C$ como el punto alcanzado por $L_1$ y $L_2$. Con $\beta=q_2+q_3$:
\begin{align}
\rho&=L_1\sin q_2+L_2\sin\beta,\\
z_C&=h+L_1\cos q_2+L_2\cos\beta,\\
\vct p_C&=\begin{bmatrix}\rho\cos q_1\\\rho\sin q_1\\z_C\end{bmatrix}.
\end{align}
Los senos aparecen en el alcance horizontal porque los eslabones parten de la dirección vertical. La orientación y posición del TCP son
\begin{align}
\mat R_{ME}&=\mat R_z(q_1)\mat R_y(\beta)
\mat R_z(q_4)\mat R_y(q_5)\mat R_z(q_6),\\
\vct p_E&=\vct p_C+d\mat R_{ME}\vct e_z,
\qquad\vct e_z=\begin{bmatrix}0\\0\\1\end{bmatrix}.
\label{eq:tcp}
\end{align}
La pose en la base se obtiene como $\mat T_{BE}=\mat T_{BM}\mat T_{ME}$. Las coordenadas del TCP deben indicar siempre el marco de expresión; una misma terna numérica cambia de significado al cambiar de marco.

\subsection{Orientación explícita de la muñeca}
Si $c_i=\cos q_i$ y $s_i=\sin q_i$, para $i=4,5,6$, la secuencia ZYZ vale
\begin{equation}
\mat R_{456}=\begin{bmatrix}
c_4c_5c_6-s_4s_6&-c_4c_5s_6-s_4c_6&c_4s_5\\
s_4c_5c_6+c_4s_6&-s_4c_5s_6+c_4c_6&s_4s_5\\
-s_5c_6&s_5s_6&c_5
\end{bmatrix}.
\end{equation}
Se premultiplica por $\mat R_z(q_1)\mat R_y(\beta)$ para expresar la orientación completa. La tercera columna indica la dirección del eje distal que multiplica al offset $d$ en~\eqref{eq:tcp}.

\subsection{Límites y postura de transporte}
Los límites geométricos son $q_1\in[-90^{\circ},90^{\circ}]$, $q_2\in[-10^{\circ},125^{\circ}]$ y $q_3,q_4,q_5,q_6\in[-90^{\circ},90^{\circ}]$. La postura de transporte del modelo es
\begin{equation}
\vct q_{\mathrm{tr}}=\begin{bmatrix}0&15&-40&0&25&0\end{bmatrix}^{\trans}
\quad[^{\circ}].
\end{equation}
En ella $\beta=-25^{\circ}$ y $\beta+q_5=0$, por lo que la orientación ideal resultante es identidad. Los límites geométricos no acreditan ausencia de colisiones, margen de cables ni comportamiento bajo carga.

Para $\vct q=\vct0$, la cadena ideal queda alineada con $z$ y produce $z_E=h+L_1+L_2+d=0{,}5375\unit m$ respecto de la montura. Esa postura es singular y no debe interpretarse como una altura física medida o una postura operativa aprobada.

\section{Cinemática inversa}
Para una pose objetivo $(\vct p_d,\mat R_d)$ expresada en la montura, se resta el desplazamiento de herramienta:
\begin{equation}
\vct p_{C,d}=\vct p_d-d\mat R_d\vct e_z.
\end{equation}
Se obtienen el azimut, alcance radial y altura relativa:
\begin{equation}
q_1=\atan(y_C,x_C),\qquad
\rho=\sqrt{x_C^2+y_C^2},\qquad s=z_C-h.
\end{equation}
La ley de cosenos del subproblema de dos eslabones da
\begin{align}
c_3&=\frac{\rho^2+s^2-L_1^2-L_2^2}{2L_1L_2},\\
q_3&=\pm\arccos(c_3),\\
q_2&=\atan(\rho,s)-\atan\!\left(L_2\sin q_3,L_1+L_2\cos q_3\right).
\end{align}
Se requiere $|c_3|\le1$, salvo una tolerancia numérica explícita. Las dos ramas representan configuraciones distintas de codo. Al permitir radio firmado aparecen candidatos adicionales con un giro alternativo de base; todos deben filtrarse por límites y continuidad.

Para orientar la muñeca se calcula
\begin{equation}
\mat R_{456,d}=\bigl[\mat R_z(q_1)\mat R_y(q_2+q_3)\bigr]^{\trans}\mat R_d
=\mat R_z(q_4)\mat R_y(q_5)\mat R_z(q_6).
\end{equation}
Con elementos $r_{ij}$ de $\mat R_{456,d}$, una rama no singular de $0<q_5<\pi$ es
\begin{align}
q_5&=\atan\!\left(\sqrt{r_{13}^2+r_{23}^2},r_{33}\right),\\
q_4&=\atan(r_{23},r_{13}),\qquad
q_6=\atan(r_{32},-r_{31}).
\end{align}
La rama $(q_4+\pi,-q_5,q_6+\pi)$ permite buscar alternativas compatibles con los límites. Si $\sin q_5=0$, los giros primero y tercero quedan acoplados y se necesita un criterio de continuidad.

\subsection{Factibilidad y selección}
La solución algebraica se evalúa con límites articulares, colisiones, entorno, estabilidad y continuidad. Entre candidatos factibles puede minimizarse
\begin{equation}
C(\vct q)=\sum_{i=1}^6w_i\,\wrap(q_i-q_{i,\mathrm{actual}})^2.
\end{equation}
Los pesos reflejan desplazamientos admisibles y preferencias de postura. El costo no vuelve viable una solución que colisiona o compromete el apoyo. La función de agarre disponible en la Web resuelve una rama descendente restringida con J4 y J6 fijados; el desarrollo analítico anterior describe el modelo general y no implica que un solucionador general esté implementado en hardware.

\section{Jacobiano y singularidades}
Para cada articulación revoluta se expresan su eje unitario $\vct a_i$ y un punto $\vct o_i$ en un mismo marco. La columna del Jacobiano geométrico es
\begin{equation}
\mat J_i=\begin{bmatrix}
\vct a_i\times(\vct p_E-\vct o_i)\\\vct a_i
\end{bmatrix},
\qquad\mat J=\begin{bmatrix}\mat J_1&\cdots&\mat J_6\end{bmatrix}.
\end{equation}
El eje $\vct a_i$ es la dirección real transformada de cada articulación, incluidos los ejes locales $Y$. Con el orden lineal--angular adoptado:
\begin{equation}
\begin{bmatrix}\vct v_E\\\vct\omega_E\end{bmatrix}=\mat J(\vct q)\dot{\vct q},
\qquad \vct\tau_{\mathrm{ext}}=\mat J^{\trans}
\begin{bmatrix}\vct F_E\\\vct M_E\end{bmatrix}.
\end{equation}
La segunda igualdad resulta de trabajo virtual y no incluye automáticamente peso propio, rozamiento ni aceleración.

\subsection{Jacobiano de posición del centro de muñeca}
Definiendo $a=L_1\cos q_2+L_2\cos\beta$ y $b=L_2\cos\beta$:
\begin{equation}
\mat J_C=\begin{bmatrix}
-\rho\sin q_1&a\cos q_1&b\cos q_1\\
\rho\cos q_1&a\sin q_1&b\sin q_1\\
0&-\rho&-L_2\sin\beta
\end{bmatrix}.
\end{equation}
De su determinante se obtiene
\begin{equation}
|\det\mat J_C|=|\rho L_1L_2\sin q_3|.
\end{equation}
El subproblema de posición pierde rango cuando $\rho=0$ o $\sin q_3=0$. La muñeca agrega singularidades con $\sin q_5=0$. Cerca de ellas, una pequeña velocidad cartesiana puede exigir velocidades articulares elevadas.

Una opción de resolución amortiguada, para una tarea correctamente escalada, es
\begin{equation}
\dot{\vct q}=\mat J^{\trans}
\left(\mat J\mat J^{\trans}+\lambda^2\mat I\right)^{-1}\vct V_d.
\end{equation}
Las componentes lineales y angulares requieren una escala consistente antes de comparar condicionamiento o elegir $\lambda$. Esta ley representa una alternativa de control a desarrollar, sin atribuirle ejecución física en el sistema.

\section{Transmisiones y pinza}
La relación de reducción $i_r=\omega_m/\omega_q$ y eficiencia $\eta$ vincula velocidad y par:
\begin{equation}
|\omega_q|=\frac{|\omega_m|}{i_r},\qquad
|\tau_q|\simeq\eta i_r|\tau_m|.
\end{equation}
La transformación de una lectura de accionamiento al ángulo geométrico puede escribirse
\begin{equation}
q_i=\sigma_i\frac{\theta_{m,i}-\theta_{0,i}}{i_{r,i}},\qquad \sigma_i\in\{-1,+1\}.
\end{equation}
El cero, el sentido y la unidad nativa se calibran para el motor y su montaje. Un engranaje exterior invierte sentido; una correa abierta normalmente lo conserva. El cociente de dientes no determina por sí solo el signo completo de una articulación.

La pinza utiliza un piñón de 24 dientes, módulo 1 y radio primitivo $r_p=12\unit{mm}$. Las dos cremalleras se desplazan simétricamente. Para un giro relativo $\theta$ en radianes:
\begin{align}
g(\theta)&=g_0+2r_p\theta,\qquad \dot g=2r_p\dot\theta,\\
F_{\mathrm{mandíbula}}&\simeq\frac{\eta\tau_p}{2r_p}.
\end{align}
El intervalo geométrico de apertura es aproximadamente 4,007 a 79,405 mm. La fuerza anterior supone reparto simétrico entre cremalleras. Para dos contactos verticales con fricción $\mu$, la condición estática ideal de retención es
\begin{equation}
2\mu F_{\mathrm{mandíbula}}\ge m_o g_0,
\qquad FS_{\mathrm{agarre}}=\frac{2\mu F_{\mathrm{mandíbula}}}{m_og_0}.
\end{equation}
Fricción, rigidez, geometría de contacto, calentamiento y límite continuo del servo deben medirse para asignar una carga útil. El par de bloqueo no se usa como capacidad continua de agarre.
\FloatBarrier

\chapter{Manipulación móvil, dinámica y estabilidad}
\label{cap:dinamica}

\section{Modelo cinemático integrado}
La configuración de pose reúne tres coordenadas de base y seis de brazo:
\begin{equation}
\vct Q=\begin{bmatrix}x&y&\psi&q_1&q_2&q_3&q_4&q_5&q_6\end{bmatrix}^{\trans}.
\end{equation}
La apertura de pinza se trata como una coordenada auxiliar. Para un marco mundial $W$, base $B$, montura $M$ y herramienta $E$:
\begin{equation}
\mat T_{WE}(\vct Q)=\mat T_{WB}(x,y,\psi)\mat T_{BM}\mat T_{ME}(\vct q),
\label{eq:integrada}
\end{equation}
con
\begin{equation}
\mat T_{WB}=\begin{bmatrix}
\cos\psi&-\sin\psi&0&x\\\sin\psi&\cos\psi&0&y\\0&0&1&0\\0&0&0&1
\end{bmatrix}.
\end{equation}
La ecuación~\eqref{eq:integrada} expresa que desplazar la base cambia la pose mundial de toda la cadena, según la formulación de manipulación móvil~\cite{mobile}. El marco de base se ha elegido a altura constante; cualquier elevación fija adicional debe incluirse en una transformación constante y no duplicarse en la montura.

\subsection{Jacobiano del conjunto}
Sea $\vct p=[p_x,p_y,p_z]^{\trans}$ la posición del TCP respecto del origen de la base. Las velocidades de cuerpo $(v_x,v_y,\omega)$ generan
\begin{equation}
\mat J_{\mathrm{base}}=
\begin{bmatrix}
1&0&-p_y\\0&1&p_x\\0&0&0\\0&0&0\\0&0&0\\0&0&1
\end{bmatrix}.
\end{equation}
Si $\mat J_{\mathrm{brazo}}$ se expresa en ejes de la base y con velocidad lineal referida al TCP, la concatenación es
\begin{equation}
\begin{bmatrix}\vct v_E\\\vct\omega_E\end{bmatrix}_B
=\underbrace{\begin{bmatrix}\mat J_{\mathrm{base}}&\mat J_{\mathrm{brazo}}\end{bmatrix}}_{\mat J_{\mathrm{total}}\in\mathbb R^{6\times9}}
\begin{bmatrix}v_x\\v_y\\\omega\\\dot{\vct q}\end{bmatrix}.
\end{equation}
Las tres primeras entradas son velocidades de cuerpo, no directamente $\dot x,\dot y$ de mundo. Si se usa $\dot{\vct Q}$, debe introducirse la rotación planar que vincula ambas parametrizaciones.

Para expresar el Jacobiano del brazo en la base basta rotar sus bloques lineal y angular cuando ambos se refieren al mismo TCP:
\begin{equation}
\mat J_{\mathrm{brazo},B}=
\begin{bmatrix}\mat R_{BM}&\mat0\\\mat0&\mat R_{BM}\end{bmatrix}
\mat J_{\mathrm{brazo},M}.
\end{equation}
La traslación de montura ya interviene al ubicar el TCP y los ejes. Al cambiar también el punto de referencia de una velocidad espacial, se necesita la transformación adjunta correspondiente.

El conjunto puede ser redundante para una tarea de seis componentes, sujeto a rango y restricciones. Repartir movimiento entre brazo y base ofrece alternativas de alcance, pero una solución de mínima velocidad articular puede acercar el centro de masa al borde de apoyo. La coordinación física simultánea debe considerar estabilidad, obstáculos, singularidades y respuesta de actuadores.

\section{Planificación de trayectorias}
Una transición articular de reposo a reposo se puede modelar con un perfil quíntico. Para $0\le t\le T$ y $\sigma=t/T$:
\begin{equation}
q_i(t)=q_{0,i}+\Delta q_i\,s(\sigma),
\qquad s(\sigma)=10\sigma^3-15\sigma^4+6\sigma^5,
\end{equation}
donde $\Delta q_i=q_{f,i}-q_{0,i}$. Sus derivadas son
\begin{align}
\dot q_i&=\frac{\Delta q_i}{T}\left(30\sigma^2-60\sigma^3+30\sigma^4\right),\\
\ddot q_i&=\frac{\Delta q_i}{T^2}\left(60\sigma-180\sigma^2+120\sigma^3\right).
\end{align}
En los extremos la velocidad y aceleración son nulas. Las cotas de magnitud son
\begin{equation}
\max|\dot q_i|=\frac{15}{8}\frac{|\Delta q_i|}{T},
\qquad
\max|\ddot q_i|=\frac{10}{\sqrt3}\frac{|\Delta q_i|}{T^2}.
\end{equation}
Por tanto, una duración que respeta límites de velocidad y aceleración debe satisfacer
\begin{equation}
T\ge\max_i\left\{
\frac{15|\Delta q_i|}{8\dot q_{\max,i}},
\sqrt{\frac{10|\Delta q_i|}{\sqrt3\ddot q_{\max,i}}}
\right\}.
\end{equation}
También se deben respetar par, temperatura, tensión de alimentación y margen de estabilidad. Una interpolación de articulaciones no produce, en general, una línea recta del TCP. Una trayectoria cartesiana requiere interpolar pose y resolver cinemática inversa con continuidad y límites.

Una secuencia de recogida puede organizarse en aproximación de base, detención, posicionamiento de brazo, cierre de pinza, retorno a transporte y continuación. La distancia ultrasónica no demuestra que el objeto esté centrado en la pinza ni que se haya retenido. Cada fase necesita una condición de terminación observable.

\section{Masa, centros e inercias}
Para cuerpos de masa $m_i$ y centro $\vct r_i$ expresado en un marco común, el centro de masa global es
\begin{equation}
\vct r_G=\frac{\sum_i m_i\vct r_i}{\sum_i m_i}.
\label{eq:cg}
\end{equation}
Las posiciones se transforman según la cadena de cada cuerpo. La masa de un actuador comercial se cuenta por instancia y no se divide entre sus ocurrencias geométricas. El motor remoto de J3 acompaña al primer eslabón; ubicarlo en el codo alteraría el cálculo de par.

Si la masa se calcula por CAD:
\begin{equation}
m_i=\rho_iV_i.
\end{equation}
La densidad debe corresponder al material efectivo. En FDM, el volumen exterior no determina automáticamente la masa laminada ni las propiedades de rigidez. El brazo combina rellenos del 25\,\%, 80\,\% y 100\,\% según la función de cada pieza. La masa no se obtiene multiplicando toda la masa sólida por ese porcentaje: también intervienen las paredes, las capas superiores e inferiores y las zonas macizas. Para actualizar centros de masa e inercias se requieren masas por pieza del laminado o del pesaje y su distribución en el ensamble. Para un cuerpo de inercia central $\mat I_G$ y vector $\vct r$ desde un origen al centro de masa:
\begin{equation}
\mat I_O=\mat R\mat I_G\mat R^{\trans}
+m\left[\left(\vct r^{\trans}\vct r\right)\mat I_3-\vct r\vct r^{\trans}\right].
\end{equation}
La matriz central se rota al marco de cálculo antes de aplicar ejes paralelos. Distribuir masa de un servo uniformemente sobre su envolvente es una aproximación que no identifica su rotor ni reductor.

\section{Dinámica del manipulador}
Las energías cinética y potencial son
\begin{align}
K&=\frac12\sum_i\left(m_i\vct v_{G_i}^{\trans}\vct v_{G_i}
+\vct\omega_i^{\trans}\mat I_{G_i}\vct\omega_i\right),\\
U&=\sum_i m_ig_0z_{G_i},\qquad g_0=9{,}80665\unit{m/s^2}.
\end{align}
Con $\mathcal L=K-U$, el formalismo de Euler--Lagrange conduce a
\begin{equation}
\frac{\dd}{\dd t}\left(\frac{\partial\mathcal L}{\partial\dot q_i}\right)
-\frac{\partial\mathcal L}{\partial q_i}
=\tau_{\mathrm{act},i}+\tau_{\mathrm{ext},i}-\tau_{f,i}.
\end{equation}
En forma matricial:
\begin{equation}
\mat M(\vct q)\ddot{\vct q}+\mat C(\vct q,\dot{\vct q})\dot{\vct q}
+\vct g_q(\vct q)+\vct\tau_f
=\vct\tau_{\mathrm{act}}+\mat J^{\trans}\vct F_{\mathrm{ext}}.
\label{eq:dinamica}
\end{equation}
$\mat M$ es la matriz de inercia, $\mat C\dot{\vct q}$ agrupa efectos centrífugos y de Coriolis, $\vct g_q$ es compensación gravitatoria y $\vct\tau_f$ representa rozamiento. Con Jacobianos del centro de masa y rotación:
\begin{equation}
\mat M=\sum_i\left[m_i\mat J_{v,i}^{\trans}\mat J_{v,i}
+\mat J_{\omega,i}^{\trans}\mat R_i\mat I_{G_i}\mat R_i^{\trans}\mat J_{\omega,i}\right].
\end{equation}
Una elección de términos de Coriolis usa los símbolos de Christoffel:
\begin{equation}
C_{ij}=\frac12\sum_k\left(
\frac{\partial M_{ij}}{\partial q_k}+
\frac{\partial M_{ik}}{\partial q_j}-
\frac{\partial M_{jk}}{\partial q_i}\right)\dot q_k,
\qquad g_{q,i}=\frac{\partial U}{\partial q_i}.
\end{equation}
Así queda definida la estructura del modelo sin inventar parámetros de inercia o rozamiento no identificados.

\subsection{Gravedad y selección de actuadores}
Con $\vct g=[0,0,-g_0]^{\trans}$, el par generalizado que ejercen los pesos es
\begin{equation}
\vct\tau_{\mathrm{grav}}=\sum_i\mat J_{v,i}^{\trans}m_i\vct g,
\qquad \vct g_q=-\vct\tau_{\mathrm{grav}}.
\end{equation}
En equilibrio, el actuador compensa el peso y la carga externa. El par disponible continuo se evalúa a partir de tensión, velocidad, eficiencia y ciclo de trabajo. Para el STS3250, la ficha utilizada especifica par nominal de 16 kgf\textperiodcentered{}cm a 12 V y par de bloqueo de 50 kgf\textperiodcentered{}cm~\cite{servo}. La conversión es
\begin{equation}
1\unit{kgf\,cm}=0{,}0980665\unit{N\,m}.
\end{equation}
El par de bloqueo no representa un punto continuo de operación. La comparación nominal no demuestra resistencia de dientes, ejes o fijaciones.

\section{Rigidez y resistencia de los eslabones}
Para un tubo circular de diámetro exterior $D$ e interior $d_i$:
\begin{equation}
A_t=\frac{\pi}{4}(D^2-d_i^2),\qquad
I=\frac{\pi}{64}(D^4-d_i^4),\qquad
J_p=\frac{\pi}{32}(D^4-d_i^4).
\end{equation}
Las relaciones elementales de flexión, torsión y desplazamiento de una viga ideal son
\begin{equation}
\sigma_{\mathrm{flex}}=\frac{M_fc}{I},\qquad
\tau_{\mathrm{tors}}=\frac{T_tc}{J_p},\qquad
\delta=\frac{F\ell^3}{3EI}.
\end{equation}
La última expresión supone voladizo, carga puntual, pequeñas deformaciones y material lineal. El montaje real incluye abrazaderas, taladros y apoyos que cambian el problema. Las piezas impresas requieren propiedades admisibles por proceso y orientación, junto con análisis de uniones y concentraciones de tensión.

\section{Estabilidad sobre ruedas}
En condición estática y piso horizontal, la proyección del centro de masa debe quedar dentro del polígono de apoyo. Para un borde $e$ con punto $\vct a_e$ y normal interior unitaria $\vct n_e$:
\begin{equation}
h_e(\vct r_G)=\vct n_e^{\trans}\left(\vct r_{G,xy}-\vct a_e\right),
\qquad h_{\min}=\min_eh_e.
\end{equation}
Se requiere $h_{\min}>0$ y un margen que cubra incertidumbre y perturbaciones. El punto de contacto de los rodillos cambia durante la marcha; un rectángulo nominal de ejes no es necesariamente la huella real de apoyo.

La comparación de momentos respecto de un borde es
\begin{equation}
M_{\mathrm{est}}-M_{\mathrm{vol}}>0,
\qquad FS_{\mathrm{vuelco}}=\frac{M_{\mathrm{est}}}{M_{\mathrm{vol}}}.
\end{equation}
El factor de vuelco no es un factor resistente del material. En aceleración o pendiente aparecen contribuciones inerciales y cambia el centro de presión. Bajo reparto uniforme y rodadura ideal puede estimarse
\begin{equation}
F\simeq M_{\mathrm{tot}}\left(a+g_0\sin\alpha+C_{rr}g_0\cos\alpha\right),
\qquad\tau_w\simeq\frac{Fr}{4}.
\end{equation}
Esta estimación no caracteriza por sí sola tracción lateral Mecanum, frenado, carga normal por rueda ni deslizamiento.

\section{Propagación de incertidumbre}
Los errores de ángulo $\delta\vct q$ y parámetros geométricos $\delta\vct\ell$ se propagan localmente como
\begin{equation}
\delta\vct p_E\simeq\mat J_v\delta\vct q+\mat J_\ell\delta\vct\ell.
\end{equation}
Si ambos conjuntos de errores se consideran independientes:
\begin{equation}
\mat\Sigma_p\simeq\mat J_v\mat\Sigma_q\mat J_v^{\trans}
+\mat J_\ell\mat\Sigma_\ell\mat J_\ell^{\trans}.
\end{equation}
Si hay correlación se agregan los términos cruzados. Holgura, flexión, montaje, TCP y calibración pueden dominar la resolución angular del sensor. Cerca de singularidades la sensibilidad puede aumentar.

Para ultrasonido, con incertidumbres pequeñas e independientes en $c$ y $\Delta t$:
\begin{equation}
\sigma_{d_s}^2\simeq\left(\frac{\Delta t}{2}\right)^2\sigma_c^2
+\left(\frac c2\right)^2\sigma_{\Delta t}^2.
\end{equation}
Las perturbaciones por superficie y eco espurio pueden superar esa contribución. Antes de atribuir precisión se necesita una referencia externa, condiciones definidas y un registro que separe exactitud, repetibilidad y resolución.
\FloatBarrier

\chapter{Resultados y análisis del desempeño}
\label{cap:resultados}

\section{Condiciones de interpretación}
Los resultados se presentan por subsistema. Las cifras de masa y estabilidad corresponden al modelo CAD con densidades y masas comerciales o estimadas. Las mediciones de movimiento corresponden a una tanda de la base del 12 de noviembre de 2026, sin carga, con consigna de 100 mm/s y firmware 1.28.0. No se extrapolan esos resultados al manipulador cargado ni a otra superficie.

El modelo matemático determina cómo deberían relacionarse las magnitudes bajo sus hipótesis. La presencia de una ecuación, un modelo 3D o una función de software no sustituye la verificación del mecanismo físico. Esta sección permite evaluar lo obtenido y localizar las limitaciones que condicionan el cierre del proyecto.

\section{Masa y estabilidad calculadas}
El presupuesto de referencia del manipulador incluye montura y control propio. Se considera una base de 1,5 kg distribuida simétricamente y una carga objetivo de 0,2 kg. La masa de base es una hipótesis aproximada y no se confunde con el pesaje de todo el conjunto.

\begin{table}[!htbp]\centering\small
\caption{Presupuesto de masa del modelo. Las diferencias de centésimas provienen del redondeo.}
\begin{tabularx}{\linewidth}{@{}YP{32mm}@{}}
\toprule\textbf{Elemento} & \textbf{Masa calculada}\\\midrule
PLA modelado con densidad sólida & 788,97 g\\
Acero y herrajes & 566,78 g\\
Cinco STS3250 & 377,50 g\\
Dos STS3032 & 43,20 g\\
Aluminio y aluminio 2024 & 217,76 g\\
Rodamientos & 155,00 g\\
Electrónica propia & 177,00 g\\
PVC & 7,93 g\\
Correa, latón y almohadillas & 10,86 g\\\midrule
\textbf{Brazo, montura y control sin objeto} & \textbf{2344,99 g}\\
Base móvil de referencia & 1500,00 g\\
Objeto de referencia & 200,00 g\\\midrule
\textbf{Conjunto de referencia} & \textbf{4044,99 g}\\\bottomrule
\end{tabularx}
\end{table}

La masa del brazo representa volumen y asignaciones físicas de CAD, no peso de una impresión terminada. La fila de PLA usa densidad sólida y no incorpora todavía la distribución de rellenos del 25\,\% para envolventes, 80\,\% para estructura cargada y 100\,\% para transmisión. Por tanto, el presupuesto y el análisis de estabilidad siguientes conservan las hipótesis de ese modelo; deben actualizarse con masas por pieza y centros de masa del conjunto fabricado. Los servos y la electrónica se representan con propiedades aproximadas. La masa de 1,5 kg de la base se interpreta sin duplicar la electrónica propia del brazo incluida en el presupuesto.

Con un apoyo rectangular nominal de $214\times124\unit{mm}$, la postura frontal extendida produce una distancia proyectada al borde cercana a 0,85 mm y
\begin{equation}
FS_{\mathrm{frontal}}=1{,}023<1{,}2=FS_{\mathrm{criterio}}.
\end{equation}
El cálculo no satisface el criterio elegido. Una pequeña variación de masa, contacto, pendiente o aceleración puede consumir el margen. Un actuador capaz de sostener la postura no evita el vuelco del conjunto.

\figura[.90]{estabilidad.png}{Proyección calculada del centro de masa para diferentes orientaciones del brazo sobre el apoyo nominal. La posición frontal se aproxima al borde; el gráfico representa un modelo estático.}{fig:estabilidad}

El estudio de configuración incluye 43.798 estados muestreados. Una muestra finita no demuestra ausencia de problemas en todo el espacio continuo. Para una postura de comparación se obtuvieron márgenes nominales de actuador de 2,58 en J2, 2,63 en J3, 19,91 en J4 y 3,51 en J5, bajo las hipótesis del modelo. Estos valores no sustituyen una envolvente de par por velocidad ni una aceptación resistente o térmica.

\section{Mediciones de avance, lateral y giro}
La tanda comprende cinco repeticiones por tipo de movimiento. La odometría es el valor interno estimado y la medida externa representa la referencia física registrada. Se diferencia error de llegada de deriva odométrica:
\begin{equation}
e_{\mathrm{llegada}}=y_{\mathrm{real}}-y_{\mathrm{objetivo}},
\qquad e_{\mathrm{odom}}=y_{\mathrm{odom}}-y_{\mathrm{real}}.
\end{equation}
Los rangos y valores centrales siguientes describen los datos disponibles; no se les asigna una incertidumbre instrumental no determinada.

\begin{table}[!htbp]\centering\small
\caption{Resultados lineales con firmware 1.28.0. Valores en milímetros.}
\begin{tabularx}{\linewidth}{@{}P{30mm}P{20mm}YP{27mm}P{26mm}@{}}
\toprule\textbf{Movimiento} & \textbf{Meta} & \textbf{Odometría final} & \textbf{Medida real} & \textbf{Error medio de llegada}\\\midrule
Avance & $+1000$ & 989 a 991 & 989 a 991 & $-10{,}2$\\
Lateral derecho & $-1000$ & $-989$ a $-992$ & $-989$ a $-992$ & $+10{,}2$\\\bottomrule
\end{tabularx}
\end{table}

\begin{table}[!htbp]\centering\small
\caption{Resultados angulares con firmware 1.28.0. Valores en grados.}
\begin{tabularx}{\linewidth}{@{}P{29mm}P{20mm}YP{25mm}P{26mm}@{}}
\toprule\textbf{Movimiento} & \textbf{Meta} & \textbf{Odometría final} & \textbf{Medida real} & \textbf{Error de llegada}\\\midrule
CCW & $+90$ & 88,6 a 89,0 & $+91$ & $+1{,}0$\\
CW & $-90$ & $-88{,}8$ a $-89{,}2$ & $-91$ & $-1{,}0$\\\bottomrule
\end{tabularx}
\end{table}

En avance se registró una deriva media odometría--medida de aproximadamente $-0{,}2\unit{mm}$ y en lateral de 0,0 mm. En los giros, la diferencia fue aproximadamente $-2{,}2^{\circ}$ en CCW y $+2{,}0^{\circ}$ en CW. Los datos muestran que la estimación de distancia puede coincidir con la referencia y aun así terminar antes de la meta.

\subsection{Interpretación del error lineal}
Las bandas de llegada de la configuración ensayada eran 10 mm y 25 mrad. Al aproximarse a una meta de 1000 mm, el controlador admite finalizar desde 990 mm. Por eso un error cercano a 10 mm puede originarse en la condición de terminación sin exigir un cambio de escala.

\begin{equation}
|1000-990|=10\unit{mm}.
\end{equation}
Corregir el factor de odometría para eliminar ese efecto mezclaría error de calibración con decisión de llegada. Deben analizarse por separado consigna, pose estimada y medida externa.

\subsection{Interpretación del giro}
En giro se combinan la banda de parada y una subestimación angular de odometría. Una relación representativa es
\begin{equation}
\frac{|\psi_{\mathrm{real}}|}{|\psi_{\mathrm{odom}}|}
\simeq\frac{91}{88{,}9}\simeq1{,}024.
\end{equation}
El valor es cercano al factor de fábrica 1,028, obtenido de una relación de giro de $92{,}5^{\circ}/90^{\circ}$. La configuración persistente puede conservar una escala distinta de la de fábrica. Por ello, actualizar software y modificar parámetros activos no son operaciones equivalentes.

\section{Ajustes y resultados por demostrar}
La fuente 1.28.1 contempla tolerancias de 4 mm y 10 mrad, migración de factores de giro y mayor ganancia de rumbo durante traslación. Su estado es fuente ATmega editada sin compilación ni carga acreditadas para ese ajuste; ESP y Web conservan actualización de versión. La tabla experimental corresponde a 1.28.0 y no demuestra el resultado de 1.28.1.

Las pruebas sin carga no caracterizan la precisión del TCP, la fuerza de pinza, el comportamiento con objeto ni la estabilidad dinámica del manipulador móvil. La base no incorpora una IMU; el patinamiento que no se refleja en las cuentas necesita una referencia externa para observarse.

\begin{table}[!htbp]\centering\small
\caption{Estado técnico y condiciones de cierre por subsistema.}
\begin{tabularx}{\linewidth}{@{}P{29mm}YY@{}}
\toprule\textbf{Área} & \textbf{Resultado disponible} & \textbf{Condición pendiente}\\\midrule
Mecánica & Geometría v17 e interfaces nominales. & Materiales, ajustes, fabricación e inspección.\\
Brazo & Cadena, IK ideal y Jacobianos. & Ceros, bus, potencia y contraste de poses.\\
Base & Mezcla, Hall y mediciones sin carga. & Caracterización ampliada por carga y superficie.\\
Estabilidad & Modelo de masa y margen frontal insuficiente. & Pesajes, contactos reales y región operativa.\\
Control & Servicios y leyes disponibles en fuente. & Ajuste, frenado y comportamiento bajo carga.\\
Pinza & Apertura y relaciones de fuerza ideales. & Contacto, fricción, calentamiento y retención.\\
Sensores & Modelos y asignación de señales. & Error, latencia, ángulo y tasa de inválidos.\\\bottomrule
\end{tabularx}
\end{table}

\section{Métricas para ampliar la caracterización}
Para $n$ mediciones con error $e_j=y_{\mathrm{med},j}-y_{\mathrm{ref},j}$:
\begin{align}
\overline e&=\frac1n\sum_{j=1}^n e_j,\\
\RMSE&=\sqrt{\frac1n\sum_{j=1}^n e_j^2},\\
s_e&=\sqrt{\frac{1}{n-1}\sum_{j=1}^n(e_j-\overline e)^2}.
\end{align}
El sesgo describe error sistemático, RMSE combina sesgo y dispersión, y $s_e$ describe dispersión alrededor de la media. La resolución de encoder es una cuantización; la repetibilidad exige repetir una condición y la exactitud exige una referencia. Ninguna de esas métricas se atribuye al brazo sin sus mediciones.

Cada sesión debe registrar versión, calibración, batería, carga, superficie, instrumento, referencia, tiempos, cuentas crudas, consignas y causa de parada. La calibración debe cubrir avance, lateral y giro; una medida lineal no basta para estimar todos los errores de movimiento.
\FloatBarrier

\chapter{Integración, evaluación y conclusiones}
\label{cap:cierre}

\section{Aportes del desarrollo}
El proyecto integra un control AVR propio para la plataforma, conectividad ESP8266, visualización Web y una interfaz local OLED/encoder. Las definiciones comunes de comandos, unidades, referencias y estados permiten relacionar las partes del sistema sin depender de convenciones implícitas.

El diseño del brazo combina seis ejes, transmisiones y una pinza simétrica con un modelo que vincula la geometría con la pose de herramienta. La transmisión remota del codo acerca masa al hombro; la reducción de hombro incrementa par de salida a cambio de velocidad. Estas decisiones se evalúan como parte de un conjunto móvil, donde la estabilidad puede limitar el trabajo antes que la capacidad nominal del actuador.

El sensado de batería, el diagnóstico de comunicación y el cargador de mantenimiento facilitan observar condiciones de funcionamiento. La interacción local conserva coherencia de navegación y confirmación de cambios. La Web aporta exploración y simulación, manteniendo identificados los datos del equipo real y los del modelo.

\section{Matriz cualitativa de riesgos técnicos}
Para priorizar integración se utiliza severidad $S$ y probabilidad cualitativa $P$ de 1 a 5. Severidad 1 representa impacto menor y 5, una consecuencia crítica para la operación; probabilidad 1 representa baja plausibilidad y 5, una condición recurrente o muy plausible. Son juicios de diseño, sin equivaler a frecuencias medidas. La prioridad relativa se expresa como $R=S\,P$.

\begin{table}[!htbp]\centering\small
\caption{Priorización técnica del trabajo de integración.}
\begin{tabularx}{\linewidth}{@{}YP{20mm}Y@{}}
\toprule\textbf{Condición} & \textbf{$S/P/R$} & \textbf{Tratamiento}\\\midrule
Margen frontal insuficiente frente al vuelco. & 5/4/20 & Restringir posturas y revisar masa y apoyo reales.\\
Potencia y bus del manipulador sin cierre. & 4/4/16 & Definir distribución, interfaces y protecciones.\\
Deriva por patinamiento lateral. & 4/3/12 & Referencia externa y calibración por condiciones.\\
Interferencia o ajuste mecánico no liberado. & 4/3/12 & Tolerancias funcionales e inspección del montaje.\\
Cero o signo incorrecto de motor y Hall. & 4/3/12 & Identificación por rueda/eje y parámetros registrados.\\
Insuficiencia de tiempo o memoria. & 4/3/12 & Evaluar presupuesto y respuesta con la carga real.\\
Interpretar simulación como capacidad medida. & 3/3/9 & Mantener identificadas las condiciones y naturaleza del resultado.\\\bottomrule
\end{tabularx}
\end{table}

\section{Secuencia de cierre}
La integración debe comenzar con interfaces mecánicas y eléctricas definidas. A continuación se calibran ceros, sentidos, cuentas y factores de odometría. La base y el brazo se caracterizan separadamente antes de evaluar una recogida completa. Esa secuencia permite atribuir un fallo al subsistema correspondiente.

\begin{table}[!htbp]\centering\small
\caption{Hitos y resultados esperados de la integración.}
\begin{tabularx}{\linewidth}{@{}P{29mm}YY@{}}
\toprule\textbf{Hito} & \textbf{Condición} & \textbf{Resultado esperado}\\\midrule
Diseño liberable & Materiales, componentes y ajustes elegidos. & Planos y lista de materiales coherentes.\\
Potencia y bus & Tensiones, corriente, conexión y fallos definidos. & Esquema funcional y distribución dimensionada.\\
Calibración & Signos, ceros, CPR y masas registrados. & Parámetros identificados por equipo y versión.\\
Caracterización & Instrumentos y protocolo establecidos. & Mediciones con condiciones e incertidumbre.\\
Manipulación móvil & Envolvente y pasos anteriores resueltos. & Éxito de recogida, tiempo de ciclo y causas de fallo medidos.\\\bottomrule
\end{tabularx}
\end{table}

El espacio alcanzable se obtiene de geometría y límites articulares. El espacio libre de colisiones agrega restricciones de interferencia. La región operativa agrega estabilidad, par, velocidad, cableado y carga. Esas tres regiones deben representarse por separado; una esfera de radio igual a la suma de longitudes no describe una capacidad operativa segura.

\section{Conclusiones}
AutoFinalRobo2 permite estudiar la robótica como una integración entre mecánica, electrónica, control y comunicación. La base Mecanum incorpora tres velocidades de cuerpo relacionadas con cuatro ruedas por un modelo matricial explícito. La medición Hall permite estimar movimiento y regular velocidad, mientras los factores direccionales y el análisis de llegada muestran la necesidad de separar calibración, estimación y control.

El manipulador queda descrito por una cadena de seis articulaciones, una solución de posición y orientación y un Jacobiano para velocidades, fuerzas e incertidumbre. La transformación mundial del TCP integra la pose de base, montura y brazo. La formulación dinámica organiza la contribución de cada masa e inercia, en lugar de concentrar toda la carga en la punta.

Las mediciones de base identifican errores de llegada y diferencias angulares bajo condiciones sin carga. Los ajustes de software pendientes no se consideran resultados experimentales. La geometría CAD del brazo tampoco equivale a capacidad física aprobada: faltan fabricación, potencia, bus, calibración y pruebas de manipulación.

La estabilidad es la limitación decisiva del conjunto de referencia. El margen frontal calculado es inferior al criterio adoptado, por lo que no se puede afirmar trabajo cargado en todo el alcance. El cierre del proyecto debe convertir los modelos en mediciones de un montaje definido y establecer una región de operación compatible con sus restricciones.
\FloatBarrier

\appendix
\chapter{Símbolos y parámetros de consulta}
\begin{longtable}{@{}P{31mm}P{82mm}P{28mm}@{}}
\caption{Nomenclatura del informe.}\\
\toprule\textbf{Símbolo} & \textbf{Significado} & \textbf{Unidad}\\\midrule\endfirsthead
\toprule\textbf{Símbolo} & \textbf{Significado} & \textbf{Unidad}\\\midrule\endhead
$\vct q$ & Ángulos geométricos del manipulador. & rad\\
$x,y,\psi$ & Pose planar de la base. & m, m, rad\\
$v_x,v_y,\omega$ & Velocidades expresadas en el cuerpo. & m/s, rad/s\\
$u_i,\omega_{w,i}$ & Velocidad periférica y angular de rueda. & m/s, rad/s\\
$r,D,L$ & Radio, diámetro y suma de semidistancias de base. & m\\
$N_i,\Delta n_i$ & Cuentas por vuelta e incremento de cuentas. & cuenta/vuelta; cuenta\\
$h,L_1,L_2,d$ & Longitudes equivalentes de la cadena del brazo. & m\\
$\mat T_{AB}$ & Transformación que lleva coordenadas de B a A. & Mixta\\
$\mat R_{AB},\vct t_{AB}$ & Rotación y traslación entre marcos. & Adimensional; m\\
$W,B,M,E$ & Mundo, base, montura y herramienta. & Marcos\\
$\mat J$ & Jacobiano geométrico, orden lineal--angular. & Según filas\\
$\vct F,\vct M$ & Fuerza y momento externos. & N, N\textperiodcentered{}m\\
$\vct\tau$ & Pares generalizados. & N\textperiodcentered{}m\\
$m_i,\vct r_G$ & Masa de un cuerpo y centro de masa global. & kg, m\\
$\mat I_G,\mat M$ & Tensor de inercia central y matriz de inercia articular. & Según coordenadas\\
$K,U,\mathcal L$ & Energía cinética, potencial y lagrangiano. & J\\
$\mat\Sigma$ & Matriz de covarianza. & Producto de unidades\\
$k_x,k_y,k_\psi$ & Factores direccionales de odometría. & Adimensional\\
$i_r,\eta$ & Relación de reducción y eficiencia. & Adimensional\\
$g,r_p$ & Apertura de pinza y radio primitivo del piñón. & m\\
$g_0$ & Aceleración gravitatoria de referencia. & m/s$^2$\\
$R_m,L_m,K_e,K_t$ & Parámetros eléctricos y electromecánicos del motor. & Según definición\\
\bottomrule
\end{longtable}

\section*{Parámetros nominales usados}
\addcontentsline{toc}{section}{Parámetros nominales usados}
\begin{table}[!htbp]\centering
\begin{tabularx}{\linewidth}{@{}YP{50mm}@{}}\toprule
\textbf{Parámetro} & \textbf{Valor}\\\midrule
Diámetro y radio de rueda & $D=0{,}080\unit m$; $r=0{,}040\unit m$\\
Suma de semidistancias & $L=0{,}170\unit m$\\
Cuentas por vuelta, captura x1 & $N_i=370$\\
Altura equivalente de hombro & $h=0{,}209\unit m$\\
Eslabón proximal & $L_1=0{,}070\unit m$\\
Distancia equivalente a muñeca & $L_2=0{,}115\unit m$\\
Offset distal de herramienta & $d=0{,}1435\unit m$\\
Factores avance / lateral / giro & $1{,}000$ / $0{,}775$ / $1{,}028$\\
Período nominal de control & $10\unit{ms}$\\
Ganancias digitales iniciales & $K_p=128$, $K_i=8$, $K_d=0$; Q8\\\bottomrule
\end{tabularx}
\end{table}
Los valores de fábrica no sustituyen los guardados en el equipo. Las longitudes son nominales y las propiedades dinámicas y de contacto requieren caracterización.

\section*{Conversiones}
\addcontentsline{toc}{section}{Conversiones}
\begin{align}
1\unit{mm}&=10^{-3}\unit m, & 1^{\circ}&=\frac{\pi}{180}\unit{rad},\\
1\unit{rpm}&=\frac{2\pi}{60}\unit{rad/s}, &1\unit{kgf\,cm}&=0{,}0980665\unit{N\,m}.
\end{align}

\clearpage
\section*{Filminas de presentación del proyecto}
\addcontentsline{toc}{section}{Filminas de presentación del proyecto}
La presentación del proyecto se conserva en el repositorio \texttt{ProyectoRobotica2}, dentro de la carpeta \texttt{presentacion}. El archivo de las filminas es \texttt{AutoFinalRobo2\_Presentacion.pptx}.

\medskip
\noindent\href{https://github.com/kobachajchir/ProyectoRobotica2/blob/codex/monorepo-cad/presentacion/AutoFinalRobo2_Presentacion.pptx}{\textbf{Abrir las filminas en el repositorio de GitHub}}.

\medskip
\noindent\href{run:../../../../presentacion/AutoFinalRobo2_Presentacion.pptx}{\textbf{Abrir la copia local de las filminas}} desde este informe guardado en su carpeta de salida dentro del proyecto.

\medskip
Las filminas acompañan la exposición oral y reúnen arquitectura, modelos, conexionado, control, cinemática, odometría, esfuerzos e interfaces. La carpeta también contiene las fichas temáticas de apoyo y el guion de la exposición.

\clearpage
\begin{thebibliography}{9}
\addcontentsline{toc}{chapter}{Bibliografía}
\bibitem{lynch} Lynch, K. M. y Park, F. C. \emph{Modern Robotics: Mechanics, Planning, and Control}. Cambridge University Press, 2017. Material oficial de los autores: \url{https://hades.mech.northwestern.edu/index.php/Modern_Robotics}.
\bibitem{omni} Lynch, K. M. y Park, F. C. \emph{Omnidirectional Wheeled Mobile Robots}, sección 13.2 del material complementario de Modern Robotics. \url{https://modernrobotics.northwestern.edu/nu-gm-book-resource/13-2-omnidirectional-wheeled-mobile-robots-part-1-of-2/}.
\bibitem{mobile} Lynch, K. M. y Park, F. C. \emph{Mobile Manipulation}, sección 13.5 del material complementario de Modern Robotics. \url{https://modernrobotics.northwestern.edu/nu-gm-book-resource/13-5-mobile-manipulation/}.
\bibitem{servo} FEETECH. \emph{STS3250 Product Specification}. Ficha de la variante de servo utilizada para la selección nominal. \url{https://www.feetechrc.com/Data/feetechrc/upload/file/20240120/6384135881578380868917773.pdf}.
\end{thebibliography}
\endinput

\end{document}
